\begin{frame}[fragile]{Repudiation}

\begin{itemize}
	\item Repudiation is the ability to deny being the sender of a message. 
	\item For example, a message and HMAC value are sent from A to B, and afterward A claims, "I didn't send this message to B. B made this up." 
    \item We need Digital Signature (or Signing) to solve this problem. 
    \item Notice that HMAC uses the shared key to enable possible repudiation. 
\end{itemize}

\end{frame}

\begin{frame}[fragile]{}

\begin{itemize}
	\item Signing is the process of creating a digital signature of a message. 
	\item A signature is a hash of the original message, which is then encrypted with the sender's private key (not the shared key as in HMAC).
	\item The signature can be verified by the recipient using the public key of the sender. 
	\item This can guarantee the original message is authentic and unmodified.
\end{itemize}

\begin{center}
\includegraphics[width=0.4\textwidth]{topics/SE_tools2/security/pic/sign/sign.png}
\end{center} 
\end{frame}

\begin{frame}[fragile]{}

\begin{itemize}
	\item Remember that the digital signature is nothing more than an encrypted hash with a sender's private key. 
\end{itemize}

\begin{center}
\includegraphics[width=0.8\textwidth]{topics/SE_tools2/security/pic/sign/idea2.png}
\end{center} 

\end{frame}

\begin{frame}[fragile]{Creating a Signature}

\begin{itemize}
	\item We need to check the integrity of the data.
	\item So, we create a signature from a private key (line 8) and a sign that contains the original data (line 7).  
\end{itemize}

\begin{Code}{}%
\begin{lstlisting}[firstnumber=1, label=glabels, xleftmargin=0pt, basicstyle=\scriptsize]% 

const { createSign, createVerify } = require('crypto');
const { publicKey, privateKey } = require('./keypair');

const data = 'I need to sign this document.';
/// SIGN
const signer = createSign('rsa-sha256');
signer.update(data);
const siguature = signer.sign(privateKey, 'hex');
\end{lstlisting}%   
\end{Code}%

\end{frame}

\begin{frame}[fragile]{Verifying the Signature}

\begin{itemize}
	\item When we receive the data and signature, we can create a verifier (line 1) with data (line2) to verify the results (line 3). 
\end{itemize}

\begin{Code}{}%
\begin{lstlisting}[firstnumber=1, label=glabels, xleftmargin=0pt, basicstyle=\scriptsize]% 

verifier = createVerify('rsa-sha256');
verifier.update(data);
const isVerified2 = verifier.verify(publicKey, siguature, 'hex');
console.log(isVerified2);
\end{lstlisting}%   
\end{Code}%

\end{frame}